%% ============================================================
%%  COMPUTATIONAL NEUROSCIENCE PROJECT REPORT
%%  Neural Network Basics: From Single Neurons to Recurrent
%%  Circuits — A Biophysical Simulation Suite
%%
%%  Compatible with Overleaf (pdfLaTeX, XeLaTeX, or LuaLaTeX)
%%  Compile with:  pdflatex main.tex   (run twice for references)
%%                 bibtex main
%%                 pdflatex main.tex
%%                 pdflatex main.tex
%% ============================================================

\documentclass[11pt, a4paper, twocolumn]{article}

%% ── Encoding & fonts ──────────────────────────────────────
\usepackage[utf8]{inputenc}
\usepackage[T1]{fontenc}
\usepackage{lmodern}
\usepackage{microtype}           % improved justification

%% ── Page geometry ─────────────────────────────────────────
\usepackage[
  top=2.0cm, bottom=2.2cm,
  left=1.8cm, right=1.8cm,
  columnsep=0.7cm
]{geometry}

%% ── Mathematics ───────────────────────────────────────────
\usepackage{amsmath, amssymb, amsthm, mathtools}
\usepackage{bm}                  % bold math
\usepackage{physics}             % \dv, \pdv, etc.

%% ── Science & units ───────────────────────────────────────
\usepackage{siunitx}
\sisetup{per-mode=symbol, separate-uncertainty=true}

%% ── Graphics ──────────────────────────────────────────────
\usepackage{graphicx}
\usepackage[font=small, labelfont=bf, labelsep=period,
            justification=justified]{caption}
\usepackage{subcaption}
\usepackage{float}
\usepackage{wrapfig}

%% ── Tables ────────────────────────────────────────────────
\usepackage{booktabs}
\usepackage{array}
\usepackage{tabularx}
\usepackage{multirow}

%% ── Code listings ─────────────────────────────────────────
\usepackage{listings}
\usepackage{xcolor}

\definecolor{codebg}{RGB}{248,248,248}
\definecolor{codegreen}{rgb}{0.0,0.5,0.0}
\definecolor{codeblue}{rgb}{0.0,0.2,0.7}
\definecolor{codegray}{rgb}{0.5,0.5,0.5}
\definecolor{codepurple}{rgb}{0.55,0.0,0.55}

\lstdefinestyle{pythonstyle}{
  language=Python,
  backgroundcolor=\color{codebg},
  basicstyle=\ttfamily\footnotesize,
  keywordstyle=\color{codeblue}\bfseries,
  stringstyle=\color{codegreen},
  commentstyle=\color{codegray}\itshape,
  numberstyle=\tiny\color{codegray},
  breaklines=true,
  breakatwhitespace=false,
  numbers=left,
  numbersep=5pt,
  tabsize=4,
  showspaces=false,
  showstringspaces=false,
  frame=single,
  framesep=3pt,
  rulecolor=\color{gray!40},
  xleftmargin=12pt
}
\lstset{style=pythonstyle}

%% ── Hyperlinks & references ───────────────────────────────
\usepackage[
  colorlinks=true,
  linkcolor=blue!70!black,
  citecolor=green!50!black,
  urlcolor=blue!80!black
]{hyperref}
\usepackage{cleveref}            % \cref, \Cref
\usepackage{natbib}
\bibliographystyle{abbrvnat}
\setcitestyle{authoryear,open={(},close={)}}

%% ── Colours & boxes ───────────────────────────────────────
\usepackage{tcolorbox}
\tcbuselibrary{skins, breakable}

\definecolor{boxblue}{RGB}{220,235,250}
\definecolor{boxorange}{RGB}{255,243,220}
\definecolor{boxgreen}{RGB}{220,245,225}
\definecolor{boxgray}{RGB}{240,240,240}

\newtcolorbox{keyequation}[1][]{
  colback=boxblue, colframe=blue!60!black,
  fonttitle=\bfseries\small, title=Key Equation,
  breakable, #1
}
\newtcolorbox{biophysbox}[1][]{
  colback=boxgreen, colframe=green!50!black,
  fonttitle=\bfseries\small, breakable, #1
}
\newtcolorbox{notebox}[1][]{
  colback=boxorange, colframe=orange!70!black,
  fonttitle=\bfseries\small, breakable, #1
}

%% ── Headers / footers ─────────────────────────────────────
\usepackage{fancyhdr}
\pagestyle{fancy}
\fancyhf{}
\fancyhead[L]{\small\textit{Computational Neuroscience Project Report}}
\fancyhead[R]{\small\thepage}
\renewcommand{\headrulewidth}{0.4pt}

%% ── Section formatting ────────────────────────────────────
\usepackage{titlesec}
\titleformat{\section}
  {\normalfont\large\bfseries\color{blue!70!black}}
  {\thesection}{1em}{}
\titleformat{\subsection}
  {\normalfont\normalsize\bfseries}
  {\thesubsection}{1em}{}
\titleformat{\subsubsection}
  {\normalfont\small\bfseries\itshape}
  {\thesubsubsection}{1em}{}

%% ── Misc utilities ────────────────────────────────────────
\usepackage{enumitem}
\usepackage{parskip}
\setlength{\parskip}{4pt}
\usepackage{lipsum}
% \usepackage{soul}  % removed: not installed               % \hl{} highlighting
% \usepackage{doi}  % removed

%% -- Custom commands --
\newcommand{\Vm}{\ensuremath{V_{\mathrm{m}}}}
\newcommand{\Vrest}{\ensuremath{V_{\mathrm{rest}}}}
\newcommand{\Vth}{\ensuremath{V_{\mathrm{th}}}}
\newcommand{\Vreset}{\ensuremath{V_{\mathrm{reset}}}}
\newcommand{\EK}{\ensuremath{E_{\mathrm{K}}}}
\newcommand{\ENa}{\ensuremath{E_{\mathrm{Na}}}}
\newcommand{\EL}{\ensuremath{E_{\mathrm{L}}}}
\newcommand{\ECl}{\ensuremath{E_{\mathrm{Cl}}}}
\newcommand{\gNa}{\ensuremath{\bar{g}_{\mathrm{Na}}}}
\newcommand{\gK}{\ensuremath{\bar{g}_{\mathrm{K}}}}
\newcommand{\gL}{\ensuremath{g_{\mathrm{L}}}}
\newcommand{\Cm}{\ensuremath{C_{\mathrm{m}}}}
\newcommand{\taum}{\ensuremath{\tau_{\mathrm{m}}}}
\newcommand{\tauw}{\ensuremath{\tau_{w}}}
\newcommand{\dT}{\ensuremath{\Delta_{T}}}
\newcommand{\Hz}{\ensuremath{\,\mathrm{Hz}}}
\newcommand{\mV}{\ensuremath{\,\mathrm{mV}}}
\newcommand{\ms}{\ensuremath{\,\mathrm{ms}}}
\newcommand{\pA}{\ensuremath{\,\mathrm{pA}}}
\newcommand{\nS}{\ensuremath{\,\mathrm{nS}}}
\newcommand{\pF}{\ensuremath{\,\mathrm{pF}}}
\newcommand{\uAcm}{\ensuremath{\,\mu\mathrm{A/cm}^{2}}}
\newcommand{\mScm}{\ensuremath{\,\mathrm{mS/cm}^{2}}}
\newcommand{\uFcm}{\ensuremath{\,\mu\mathrm{F/cm}^{2}}}
\newcommand{\python}[1]{\texttt{#1}}

% ──────────────────────────────────────────────────────────
\begin{document}

%% ── Title block ───────────────────────────────────────────
\twocolumn[{%
\begin{@twocolumnfalse}

\begin{center}

{\LARGE\bfseries\color{blue!70!black}
  Neural Network Basics in Computational Neuroscience}\\[6pt]
{\large
  A Biophysical Simulation Suite: From Action Potentials\\
  to Recurrent E/I Networks with STDP Plasticity}\\[12pt]

{\normalsize
  \textbf{Computational Biophysics Project Report}\\[4pt]
  Python implementation: Hodgkin--Huxley $\cdot$ LIF $\cdot$ AdEx $\cdot$
  AMPA/NMDA/GABA synapses $\cdot$ STDP $\cdot$ E/I network\\[4pt]
  Full source code and database at:
  \texttt{\href{https://github.com/}{github.com/[your-username]/compneuro}}
}

\vspace{8pt}
\hrule height 0.8pt
\vspace{8pt}

\begin{minipage}{0.88\textwidth}
\small
\textbf{Abstract.}
We present a self-contained, scientifically rigorous Python simulation suite
covering the core hierarchy of computational neuroscience: biophysical
membrane dynamics, synaptic transmission, activity-dependent plasticity, and
recurrent network behaviour. Starting from the canonical Hodgkin--Huxley
(1952) conductance-based model, we implement the full quartet of
voltage-gated Na\textsuperscript{+} and K\textsuperscript{+} kinetics with
fourth-order Runge--Kutta integration (\SI{0.025}{\milli\second} time-step),
and characterise firing via the frequency--current (f--I) relationship and
V--n phase-plane analysis. We then contrast two reduced spiking neuron
families --- the Leaky Integrate-and-Fire (LIF) and the Adaptive Exponential
Integrate-and-Fire (AdEx, \citealt{Brette2005}) --- demonstrating five
cortical firing-pattern archetypes (RS, IB, CH, FS, LTS). Synaptic
transmission is modelled through kinetic receptor schemes for AMPA, NMDA
(\citealt{Destexhe1994}), GABA\textsubscript{A}, and GABA\textsubscript{B},
with the voltage-dependent Mg\textsuperscript{2+} block of NMDA
(\citealt{Jahr1990}) explicitly reproduced. Synaptic plasticity is implemented
as asymmetric spike-timing-dependent plasticity (STDP; \citealt{Bi1998}), and
the steady-state weight distribution under competitive Hebbian learning is
characterised across 50 independent synapses. Finally, a 125-neuron recurrent
excitatory--inhibitory (E/I) LIF network is simulated with random sparse
connectivity ($p = 0.15$), yielding biologically plausible population dynamics
including spontaneous firing (E: $\sim$21\Hz, I: $\sim$26\Hz), LFP proxy
signals, and frequency-band decomposition. All simulation results are
persisted to a SQLAlchemy-managed SQLite database for reproducibility.
Ten publication-quality figures plus a summary dashboard are produced, each
referencing the primary experimental and theoretical literature.
\end{minipage}

\vspace{8pt}
\hrule height 0.8pt
\vspace{12pt}
\end{center}
\end{@twocolumnfalse}
}]

%% ──────────────────────────────────────────────────────────
\section{Introduction and Scientific Motivation}
\label{sec:intro}
%% ──────────────────────────────────────────────────────────

The question of how the nervous system encodes, transmits, and transforms
information is among the most profound in modern science. Despite the apparent
simplicity of a single neuron --- a cell roughly \SIrange{10}{100}{\micro\metre}
in diameter --- the collective behaviour of billions of such cells gives rise
to sensation, memory, decision-making, and consciousness. \emph{Computational
neuroscience} bridges this gap by constructing mathematical and numerical
models at each level of neural organisation, from ion-channel kinetics to
whole-brain dynamics \citep{Dayan2001}.

The present work addresses three interlinked levels:
\begin{enumerate}[leftmargin=1.2em, itemsep=2pt]
  \item \textbf{Single-neuron biophysics.} How do ion channels, described by
    voltage-dependent kinetics, generate the stereotyped action potential
    (spike) that is the fundamental unit of neural communication?
  \item \textbf{Synaptic transmission and plasticity.} How do chemical
    synapses implement specific temporal filter properties, and how does
    activity-dependent modification of synaptic strength --- plasticity ---
    implement a form of learning?
  \item \textbf{Network dynamics.} How does the balance of excitation and
    inhibition shape population activity, oscillations, and information
    transmission in recurrent circuits?
\end{enumerate}

Understanding these levels is not merely academic. Dysfunctions in ion-channel
kinetics underlie epilepsy and Long-QT syndrome \citep{Bhatt2009}; pathological
NMDA hypofunction is implicated in schizophrenia \citep{Coyle2006}; disrupted
E/I balance is a core hypothesis of autism spectrum disorder
\citep{Yizhar2011}; and aberrant synaptic plasticity is central to both
addiction and Alzheimer's disease \citep{Bhatt2009}. A simulation suite that
reproduces the key biophysical phenomena at each level therefore constitutes
an essential research and educational infrastructure for both basic and
translational neuroscience.

%% ──────────────────────────────────────────────────────────
\section{Background: Biophysical Principles}
\label{sec:background}
%% ──────────────────────────────────────────────────────────

\subsection{The Electrochemical Basis of the Resting Potential}

A neuron at rest maintains a membrane potential $\Vrest \approx -70\mV$
sustained by the unequal distribution of ions across the lipid bilayer:
high intracellular K\textsuperscript{+} and low Na\textsuperscript{+}
relative to the extracellular milieu. The Nernst equation gives the
equilibrium (reversal) potential for each ion species $X$:
\begin{equation}
  E_X = \frac{RT}{zF} \ln \frac{[X]_{\mathrm{out}}}{[X]_{\mathrm{in}}}
  \label{eq:nernst}
\end{equation}
where $R$ is the gas constant, $T$ temperature, $z$ valence, and $F$ the
Faraday constant. For mammalian neurons: $\EK \approx -77\mV$,
$\ENa \approx +55\mV$, and $\ECl \approx -70\mV$ under physiological conditions.
The resting membrane potential is well described by the Goldman--Hodgkin--Katz
(GHK) equation, which weights each ion's contribution by its membrane
permeability \citep{Goldman1943}.

\subsection{The Action Potential as an Electrical Instability}

When a depolarising stimulus raises $\Vm$ above a threshold
$\Vth \approx -55\mV$, a rapid positive-feedback cycle unfolds:
depolarisation opens voltage-gated Na\textsuperscript{+} channels (VGSC),
Na\textsuperscript{+} influx further depolarises the cell, opening more
channels. This regenerative event --- the action potential (AP) --- peaks near
$+30\mV$ before Na\textsuperscript{+} channel inactivation and K\textsuperscript{+}
channel activation restore the resting potential. The AP occupies
$\sim$\SI{1}{\milli\second} and is propagated along the axon by local circuit
currents, serving as the elementary unit of long-range neural communication.

\subsection{Chemical Synaptic Transmission}

Action potentials arriving at presynaptic terminals trigger
Ca\textsuperscript{2+}-dependent vesicle fusion and neurotransmitter release.
Glutamate (the principal excitatory transmitter) activates postsynaptic AMPA
and NMDA receptors; GABA (the principal inhibitory transmitter) activates
GABA\textsubscript{A} and GABA\textsubscript{B} receptors. Each receptor
type confers a distinctive time course and ionic selectivity to the
postsynaptic conductance change, encoding temporal information across
timescales from $\sim$\SI{2}{\ms} (AMPA) to $\sim$\SI{150}{\ms} (NMDA) and
$\sim$\SI{200}{\ms} (GABA\textsubscript{B}).

\subsection{Hebbian Plasticity and STDP}

The synaptic efficacy --- encoded by the peak conductance $g_{\max}$ --- is
not fixed but adapts as a function of correlated pre- and postsynaptic
activity. Hebbian learning (\textit{``neurons that fire together, wire
together''}) is the conceptual foundation. Its precise form in biological
synapses is spike-timing-dependent plasticity (STDP): synaptic potentiation
(LTP) occurs when a presynaptic spike precedes the postsynaptic AP by
$\Delta t > 0$, and depression (LTD) occurs for $\Delta t < 0$
\citep{Bi1998}. This asymmetric Hebbian rule provides a substrate for
temporal sequence learning and competitive synaptic development.

\subsection{E/I Balance and Network States}

Cortical circuits contain approximately 80\% excitatory (glutamatergic)
pyramidal cells and 20\% inhibitory (GABAergic) interneurons. Their
recurrent interactions generate an intrinsic tension: excitation promotes
network activity while inhibition stabilises it. The balance of these forces
determines the \emph{network state}: synchronous--regular (SR), asynchronous--irregular
(AI), or synchronous--irregular (SI) \citep{Brunel2000}. The AI state is
considered the biologically realistic operating regime of awake cortex, where
individual neurons fire irregularly (CV-ISI $\approx 1$) even though
population activity is stable.

%% ──────────────────────────────────────────────────────────
\section{Mathematical Models and Implementation}
\label{sec:methods}
%% ──────────────────────────────────────────────────────────

\subsection{Hodgkin--Huxley Neuron Model}
\label{sec:hh}

\subsubsection{Governing Equations}

The Hodgkin--Huxley (HH) model \citep{HodgkinHuxley1952} describes the
neuronal membrane as an electrical circuit with a capacitance $\Cm$
(\SI{1}{\uFcm}) in parallel with three conductance branches:

\begin{keyequation}
\begin{align}
  \Cm \dv{V}{t} &= I_{\mathrm{ext}}
    - \underbrace{\gNa m^3 h (V - \ENa)}_{I_{\mathrm{Na}}}
    - \underbrace{\gK  n^4   (V - \EK)}_{I_{\mathrm{K}}}
    - \underbrace{\gL       (V - \EL)}_{I_{\mathrm{L}}}
    \label{eq:hh_V}
\end{align}
\end{keyequation}

where $\gNa = 120\mScm$, $\gK = 36\mScm$, $\gL = 0.3\mScm$ are the maximum
conductances, and $\ENa = +55\mV$, $\EK = -77\mV$, $\EL = -54.4\mV$ are
reversal potentials measured in squid giant axon at \SI{6.3}{\celsius}.

The gating variables $m, h, n \in [0,1]$ each obey first-order kinetics:
\begin{equation}
  \dv{x}{t} = \alpha_x(V)(1-x) - \beta_x(V)\,x, \quad x \in \{m,h,n\}
  \label{eq:gate_kinetics}
\end{equation}

The rate functions $\alpha_x(V)$ and $\beta_x(V)$ are empirically determined
Boltzmann-type expressions (Table~\ref{tab:hh_rates}). At steady state:
$x_\infty(V) = \alpha_x / (\alpha_x + \beta_x)$ and
$\tau_x(V) = 1/(\alpha_x + \beta_x)$.

\begin{table}[H]
\centering
\small
\caption{Hodgkin--Huxley rate functions (original 1952 parameterisation).
  Voltage in mV, rates in ms$^{-1}$, conductances in mS/cm$^2$.}
\label{tab:hh_rates}
\begin{tabularx}{\columnwidth}{l X X}
\toprule
\textbf{Gate} & \boldmath$\alpha_x(V)$ & \boldmath$\beta_x(V)$ \\
\midrule
$m$ & $\dfrac{0.1(V+40)}{1-e^{-(V+40)/10}}$
    & $4\,e^{-(V+65)/18}$ \\[8pt]
$h$ & $0.07\,e^{-(V+65)/20}$
    & $\dfrac{1}{1+e^{-(V+35)/10}}$ \\[8pt]
$n$ & $\dfrac{0.01(V+55)}{1-e^{-(V+55)/10}}$
    & $0.125\,e^{-(V+65)/80}$ \\
\bottomrule
\end{tabularx}
\end{table}

\subsubsection{Biophysical Interpretation}

The $m^3 h$ factor in $I_{\mathrm{Na}}$ reflects the tetrameric structure of
the Na\textsuperscript{+} channel: three activation subunits ($m$) and one
inactivation subunit ($h$). Because $\tau_m \ll \tau_h \approx \tau_n$,
activation is near-instantaneous relative to inactivation and K\textsuperscript{+}
activation. This separation of timescales is mechanistically responsible for
the AP waveform: during the upstroke, $m$ activates rapidly while $h$ and $n$
lag; as $h$ closes and $n$ opens, repolarisation occurs. The \emph{absolute
refractory period} corresponds to the interval while $h$ remains
inactivated.

\subsubsection{Numerical Integration}

We integrate \cref{eq:hh_V,eq:gate_kinetics} simultaneously using the
classical fourth-order Runge--Kutta (RK4) scheme with
$\Delta t = 0.025\ms$:
\begin{align}
  \mathbf{y}_{n+1} &= \mathbf{y}_n + \frac{\Delta t}{6}(k_1 + 2k_2 + 2k_3 + k_4)
  \nonumber\\
  k_i &= f\!\left(\mathbf{y}_n + \Delta t \cdot \mathbf{c}_i\right)
  \label{eq:rk4}
\end{align}
where $\mathbf{y} = [V, m, h, n]^\top$. RK4 provides fourth-order accuracy
($\mathcal{O}(\Delta t^4)$) with three extra function evaluations per step,
critical for the stiff exponential dynamics near $V_{\mathrm{th}}$.

\subsection{Leaky Integrate-and-Fire Neuron}
\label{sec:lif}

The LIF model \citep{Lapicque1907} represents the subthreshold membrane as
an RC circuit, with spikes emitted when a hard voltage threshold is reached:
\begin{equation}
  \Cm \dv{V}{t} = -\gL(V - \EL) + I_{\mathrm{ext}}
  \label{eq:lif}
\end{equation}
with spike condition: if $V \geq \Vth$ then emit spike, $V \leftarrow \Vreset$,
and apply absolute refractory period $\tau_{\mathrm{ref}}$.

The membrane time constant $\taum = \Cm/\gL$ sets the integration window.
For our canonical parameters ($\Cm = 200\pF$, $\gL = 10\nS$): $\taum = 20\ms$.
The analytical solution for constant input gives the interspike interval
(ISI) and thus firing rate:
\begin{equation}
  f(I) = \left[\tau_{\mathrm{ref}} - \taum \ln\!\left(1 -
    \frac{\gL(\Vth - \EL)}{I - \gL(\Vreset - \EL)}\right)\right]^{-1}
  \label{eq:lif_rate}
\end{equation}

The LIF f--I curve is therefore monotonically increasing above the rheobase
$I_{\mathrm{rh}} = \gL(\Vth - \EL)$. Its linearity at suprathreshold currents
(compared to the sub-linear saturation of HH) reflects the absence of sodium
channel inactivation.

\subsection{Adaptive Exponential Integrate-and-Fire (AdEx)}
\label{sec:adex}

The AdEx model \citep{Brette2005} augments the EIF membrane equation
\citep{Fourcaud2003} with a subthreshold adaptation current $w$:
\begin{keyequation}
\begin{align}
  \Cm \dv{V}{t} &= -\gL(V - \EL) + \gL \dT e^{(V-V_T)/\dT} - w + I_{\mathrm{ext}}
  \label{eq:adex_V}\\[4pt]
  \tauw \dv{w}{t} &= a(V - \EL) - w
  \label{eq:adex_w}
\end{align}
\end{keyequation}
with the spike rule: if $V \geq \Vth$ then $V \leftarrow \Vreset$,
$w \leftarrow w + b$.

The exponential term $\gL \dT e^{(V-V_T)/\dT}$ reproduces the sharp
spike-initiation dynamics observed in cortical neurons
\citep{Fourcaud2003}, parameterised by the spike-initiation threshold $V_T$
and sharpness $\dT$. The adaptation parameters $(a, b, \tauw)$ determine
firing pattern:
\begin{itemize}[itemsep=1pt, leftmargin=1.4em]
  \item \textbf{RS (Regular Spiking):} spike-frequency adaptation due to
    large $b$ and long $\tauw$.
  \item \textbf{IB (Intrinsic Bursting):} burst of spikes followed by
    quiescence, due to shorter $\tauw$.
  \item \textbf{CH (Chattering):} high-frequency bursts.
  \item \textbf{FS (Fast Spiking):} negligible adaptation ($a = b = 0$);
    characteristic of PV\textsuperscript{+} interneurons.
  \item \textbf{LTS (Low-Threshold Spiking):} strong subthreshold coupling
    ($a = 8\nS$, $b = 200\pA$); characteristic of SST\textsuperscript{+}
    interneurons.
\end{itemize}

Table~\ref{tab:adex_presets} lists the canonical parameter sets from
\citet{Naud2008}.

\begin{table}[H]
\centering
\small
\caption{AdEx parameter presets for cortical cell types \citep{Naud2008}.
  Common parameters: $\Cm = 200\pF$, $\gL = 10\nS$, $\EL = -70\mV$,
  $V_T = -50\mV$, $\dT = 2\mV$, $\Vreset = -58\mV$.}
\label{tab:adex_presets}
\begin{tabular}{lrrrr}
\toprule
\textbf{Type} & \boldmath$a$ (nS) & \boldmath$b$ (pA) & \boldmath$\tauw$ (ms) & \textbf{CV-ISI} \\
\midrule
RS  & 4.0 &  80.5 & 144 & $>0$ \\
IB  & 4.0 &  80.5 &  16 & $\sim$0.05 \\
CH  & 4.0 &  80.5 &   5 & $\sim$0.006 \\
FS  & 0.0 &   0.0 &  40 & $\approx 0$ \\
LTS & 8.0 & 200.0 & 200 & $>0$ \\
\bottomrule
\end{tabular}
\end{table}

\subsection{Synaptic Models}
\label{sec:synapses}

\subsubsection{Kinetic Receptor Scheme}

All fast ionotropic synapses (AMPA, NMDA, GABA\textsubscript{A}) are modelled
using the two-state kinetic scheme of \citet{Destexhe1994}:
\begin{equation}
  \dv{r}{t} = \alpha [T](1-r) - \beta r
  \label{eq:kinetic_synapse}
\end{equation}
where $r \in [0,1]$ is the fraction of open channels, $[T]$ is the
transmitter concentration (modelled as a square pulse $[T]_{\max} = 1\,\mathrm{mM}$
for duration $\tau_{\mathrm{pulse}}$ after each presynaptic spike), and
$\alpha$, $\beta$ are binding and unbinding rates. The synaptic current is:
\begin{equation}
  I_{\mathrm{syn}} = g_{\max} \cdot r(t) \cdot (V_{\mathrm{post}} - E_{\mathrm{rev}})
  \label{eq:syn_current}
\end{equation}

AMPA parameters: $\alpha = 1.1\,(\mathrm{mM \cdot ms})^{-1}$,
$\beta = 0.19\,\mathrm{ms}^{-1}$, $E_{\mathrm{rev}} = 0\mV$,
$\tau_{\mathrm{decay}} \approx 5\ms$.
GABA\textsubscript{A} parameters: $\alpha = 5.0$, $\beta = 0.18\,\mathrm{ms}^{-1}$,
$E_{\mathrm{rev}} = -70\mV$, $\tau_{\mathrm{decay}} \approx 5\ms$.

\subsubsection{NMDA Receptor: Voltage-Dependent Mg\textsuperscript{2+} Block}

NMDA receptors carry an additional, critical non-linearity: at resting
membrane potentials, the channel pore is physically occluded by
extracellular Mg\textsuperscript{2+} ions \citep{Nowak1984}. This block is
relieved by postsynaptic depolarisation, making NMDA receptors function as
\emph{coincidence detectors} that respond only when pre- and postsynaptic
activity co-occur:
\begin{equation}
  B(V) = \frac{1}{1 + \exp(-0.062\,V) \cdot [\mathrm{Mg}^{2+}]_o / 3.57}
  \label{eq:mg_block}
\end{equation}
\citep{Jahr1990}. The full NMDA conductance is:
$g_{\mathrm{NMDA}}(t,V) = g_{\max}\,r(t)\,B(V)$.
NMDA kinetics are significantly slower than AMPA ($\tau_{\mathrm{decay}} \approx 150\ms$).

\subsubsection{GABA\textsubscript{B}: Metabotropic G-Protein Cascade}

GABA\textsubscript{B} receptors are G-protein-coupled and activate inwardly
rectifying K\textsuperscript{+} channels via a second-messenger cascade
\citep{Destexhe1994}. We implement the four-state model:
\begin{align}
  \dv{R}{t} &= K_1[T](1-R) - K_2 R \label{eq:gabab_R}\\
  \dv{G}{t} &= K_3 R - K_4 G \label{eq:gabab_G}\\
  g_{\mathrm{GABA_B}} &= g_{\max} \frac{G^n}{G^n + K_d} \label{eq:gabab_g}
\end{align}
where $R$ is receptor activation, $G$ is G-protein activation, and the
Hill function (\cref{eq:gabab_g}) with cooperativity $n = 4$ reflects
G-protein stoichiometry. The time-to-peak is $\sim$\SI{120}{\ms}, and
$E_{\mathrm{rev}} = -95\mV$ (K\textsuperscript{+} Nernst potential).

\subsection{Spike-Timing Dependent Plasticity}
\label{sec:stdp}

The classical asymmetric STDP rule \citep{Bi1998,Song2000} updates synaptic
weight $w \in [w_{\min}, w_{\max}]$ based on relative spike timing
$\Delta t = t_{\mathrm{post}} - t_{\mathrm{pre}}$:
\begin{keyequation}
\begin{equation}
  \Delta w = \begin{cases}
    A_+ \left(\dfrac{w_{\max} - w}{w_{\max} - w_{\min}}\right)^\mu
      e^{-\Delta t / \tau_+}
      & \Delta t > 0 \text{ (LTP)} \\[8pt]
    -A_- \left(\dfrac{w - w_{\min}}{w_{\max} - w_{\min}}\right)^\mu
      e^{\Delta t / \tau_-}
      & \Delta t < 0 \text{ (LTD)}
  \end{cases}
  \label{eq:stdp}
\end{equation}
\end{keyequation}

The multiplicative soft weight-bound ($\mu = 1$) ensures weights remain
bounded and promotes a unimodal weight distribution, unlike the additive
rule \citep{Song2000}. We use: $A_+ = 0.010$, $A_- = 0.0105$
($A_-/A_+ = 1.05 > 1$ ensures net LTD at equal pre/post rates,
preventing runaway potentiation), $\tau_+ = \tau_- = 20\ms$.

Traces are maintained for pre- ($x_{\mathrm{pre}}$) and post-synaptic
($x_{\mathrm{post}}$) spikes, decaying exponentially:
$\dot{x} = -x/\tau + \delta(t - t_{\mathrm{spike}})$.

\subsection{Recurrent E/I Network}
\label{sec:network}

We simulate a random sparse Erdős--Rényi network of $N_E = 100$ excitatory
and $N_I = 25$ inhibitory LIF neurons, connection probability $p = 0.15$
(no autapses), based on the Brunel (2000) sparse random network framework.
Synaptic connections are:
\begin{itemize}[itemsep=1pt, leftmargin=1.4em]
  \item E$\to$E, E$\to$I: AMPA-like, $g_{\mathrm{EE}} = g_{\mathrm{EI}} = 0.35\nS$
  \item I$\to$E, I$\to$I: GABA\textsubscript{A}-like, $g_{\mathrm{IE}} = g_{\mathrm{II}} = 2.0\nS$
\end{itemize}
An external Poisson drive at \SI{15}{\kilo\Hz} with $g_{\mathrm{ext}} = 0.9\nS$
provides tonic excitation, representing input from thousands of upstream neurons
not explicitly simulated. A DC baseline of $I_{\mathrm{DC}} = 200\pA$ is also
applied to each neuron.

\subsection{Spike Train Analysis}
\label{sec:analysis}

For each simulated spike train we compute:
\begin{itemize}[itemsep=1pt, leftmargin=1.4em]
  \item \textbf{CV-ISI}: coefficient of variation of inter-spike intervals
    $\mathrm{CV} = \sigma_{\mathrm{ISI}}/\mu_{\mathrm{ISI}}$; Poisson process $\Rightarrow$ CV $= 1$.
  \item \textbf{Fano factor}: variance/mean of spike counts in non-overlapping
    100-ms bins; sub-Poisson $\Rightarrow F < 1$ (regular), super-Poisson
    $\Rightarrow F > 1$ (bursty).
  \item \textbf{LFP proxy}: mean E-cell membrane voltage, smoothed with a
    10-ms Gaussian kernel, serves as a proxy for the local field potential.
  \item \textbf{PSD}: Welch's method \citep{Welch1967} applied to the LFP
    trace with \SI{1024}{sample} segments, yielding power in the
    $\delta/\theta/\alpha/\beta/\gamma$ bands.
\end{itemize}

%% ──────────────────────────────────────────────────────────
\section{Results and Figure Analysis}
\label{sec:results}
%% ──────────────────────────────────────────────────────────

\subsection{Figure 1 --- HH Action Potential Anatomy}

\begin{figure}[H]
  \centering
  \includegraphics[width=\columnwidth]{fig_01_hh_action_potential.png}
  \caption{\textbf{Hodgkin--Huxley single-neuron simulation.}
    \textbf{(A)} Membrane voltage $V(t)$ under a step current of
    $10\uAcm$ (onset at $t=10\ms$). Grey dashed lines denote Nernst potentials
    $\ENa$, $\EK$, $\EL$. Repeated AP trains at $\sim$67\Hz are produced.
    \textbf{(B--D)} Time course of gating variables $m(t)$, $h(t)$, $n(t)$.
    Note that $m$ activates on a sub-millisecond timescale, $h$ inactivates
    slowly, and $n$ activates with a delay --- the kinetic basis of the
    refractory period. \textbf{(E)} Ionic currents. The inward $I_{\mathrm{Na}}$
    (negative, depolarising) and outward $I_{\mathrm{K}}$ (positive,
    repolarising) are clearly delineated. \textbf{(F)} Injected current
    waveform.}
  \label{fig:hh_ap}
\end{figure}

\Cref{fig:hh_ap} validates the implementation against the classical
Hodgkin--Huxley prediction. The AP peak of $\sim$$+35\mV$ and undershoot
to $\sim$$-75\mV$ match the squid axon data of \citet{HodgkinHuxley1952}.
The asymmetry between the $m^3h$ Na\textsuperscript{+} activation (fast,
large inward current) and $n^4$ K\textsuperscript{+} activation (slower,
smaller outward current) is clearly visible in panel (E), explaining the
pronounced after-hyperpolarisation (AHP) that follows each spike.

\begin{biophysbox}[title=Why is the Na$^+$ current transient but K$^+$ sustained?]
The Na\textsuperscript{+} channel has \emph{two gates}: the activation gate
$m$ (fast, opens on depolarisation) and the inactivation gate $h$ (slow,
closes on depolarisation). Their product $m^3h$ therefore rises quickly then
falls as $h$ closes, yielding a transient inward current. The K\textsuperscript{+}
channel has only one gate $n$ that activates slowly, producing a sustained
outward current responsible for repolarisation.
\end{biophysbox}

\subsection{Figure 2 --- Phase Plane Analysis}

\begin{figure}[H]
  \centering
  \includegraphics[width=\columnwidth]{fig_02_phase_plane.png}
  \caption{\textbf{HH phase-plane analysis.}
    \textbf{(A)} V--n phase plane: the V-nullcline ($\dot{V}=0$, blue) and
    n-nullcline ($\dot{n}=0$, orange) intersect at the unstable fixed point
    near $V = -65\mV$, $n \approx 0.32$. The gray trajectory shows the
    limit cycle traversed during repetitive firing. The large excursion in
    $n$ during each spike corresponds to K\textsuperscript{+} channel
    opening. \textbf{(B)} Voltage-dependent steady-state gating curves
    $m_\infty(V)$, $h_\infty(V)$, $n_\infty(V)$. The half-activation
    voltage and slope factor of each curve determine threshold and kinetics.}
  \label{fig:phase_plane}
\end{figure}

Phase-plane analysis (\Cref{fig:phase_plane}A) reveals the dynamical
topology of repetitive firing. The limit cycle is the geometric object
corresponding to sustained oscillation: the trajectory circles around the
unstable equilibrium point, visited periodically. The widths of the
nullclines near their intersection predict the bifurcation type (here, a
subcritical Hopf bifurcation at the rheobase), which determines whether
spiking onset is continuous or discontinuous in the f--I curve.

\subsection{Figure 3 --- Frequency--Current (f-I) Curves}

\begin{figure}[H]
  \centering
  \includegraphics[width=\columnwidth]{fig_03_fI_curves.png}
  \caption{\textbf{f-I curves for three neuron models.}
    \textbf{(A)} HH model: rheobase $\approx 2.1\uAcm$; firing rate saturates
    near $\sim$100\Hz due to sodium channel inactivation (relative refractory
    period). \textbf{(B)} LIF vs AdEx-RS vs AdEx-FS: the LIF f-I curve
    is linear above rheobase (no spike-rate adaptation); AdEx-RS shows
    pronounced sub-linearity due to the adaptation current $w$ accumulating
    with spikes; AdEx-FS (no adaptation) tracks LIF closely but with the
    exponential spike initiation producing a sharper threshold.}
  \label{fig:fi_curves}
\end{figure}

The f-I curve (\Cref{fig:fi_curves}) is the primary functional characterisation
of a spiking neuron: it maps input current to output firing rate, i.e., it
performs \emph{rate coding}. The key observation is that adaptation (AdEx-RS)
compresses the dynamic range, acting as an automatic gain control mechanism.
This is biologically important: it prevents saturation under strong sustained
inputs and maintains sensitivity to changes in input rather than absolute
magnitude --- a form of contrast enhancement \citep{Benda2003}.

\subsection{Figure 4 --- AdEx Firing Pattern Gallery}

\begin{figure}[H]
  \centering
  \includegraphics[width=\columnwidth]{fig_04_adex_patterns.png}
  \caption{\textbf{AdEx firing pattern gallery.}
    Left panels: membrane voltage $V(t)$ (coloured) and adaptation variable
    $w(t)$ (dashed purple) for each cell type. Right panels: ISI histograms.
    RS shows strong spike-frequency adaptation (the ISI histogram is right-skewed,
    CV $= 0.33$); IB shows tight burst clustering (CV $= 0.04$); CH shows highly
    regular chattering (CV $= 0.006$, ISI histogram is narrowly peaked); FS shows
    perfectly regular firing (CV $\approx 0$); LTS shows modest adaptation.
    }
  \label{fig:adex_patterns}
\end{figure}

The AdEx parameter space replicates an impressive diversity of cortical
firing patterns (\Cref{fig:adex_patterns}). This is clinically relevant:
disruption of the fast-spiking (FS) pattern, characteristic of
parvalbumin-positive (PV\textsuperscript{+}) interneurons, is a leading
candidate mechanism for gamma oscillation deficits in schizophrenia
\citep{Gonzalez2011}. The intrinsic bursting (IB) pattern, found in layer
V pyramidal neurons, is implicated in epileptiform activity when
insufficiently balanced by inhibition.

\subsection{Figure 5 --- Synaptic Conductance Kinetics}

\begin{figure}[H]
  \centering
  \includegraphics[width=\columnwidth]{fig_05_synapse_kinetics.png}
  \caption{\textbf{Conductance kinetics of four receptor types.}
    \textbf{AMPA}: fast rise ($\sim$1\ms), decay $\tau \approx 5\ms$,
    peak $g \approx 666\,\mathrm{pS}$; \textbf{NMDA}: slow kinetics,
    $\tau_{\mathrm{decay}} \approx 150\ms$; \textbf{GABA-A}: fast
    chloride-mediated inhibition, peak $\approx 965\,\mathrm{pS}$;
    \textbf{GABA-B}: slow metabotropic kinetics, peak at $\sim$120\ms
    following a 4-spike burst. Note the four-order-of-magnitude timescale
    difference between AMPA and GABA\textsubscript{B}.}
  \label{fig:synapses}
\end{figure}

The timescale diversity of synaptic conductances (\Cref{fig:synapses})
enables the nervous system to perform temporal integration across multiple
scales simultaneously. AMPA receptors mediate rapid, point-like coincidence
signals; NMDA receptors integrate activity over $\sim$150\ms, implementing
a sliding temporal window for associative plasticity; GABA\textsubscript{B}
provides a delayed, prolonged inhibitory shunt that suppresses bursting.

\subsection{Figure 6 --- NMDA Mg\textsuperscript{2+} Block}

\begin{figure}[H]
  \centering
  \includegraphics[width=\columnwidth]{fig_06_nmda_mg_block.png}
  \caption{\textbf{NMDA voltage-dependent Mg\textsuperscript{2+} block.}
    \textbf{(A)} Unblock factor $B(V)$ as a function of membrane potential for
    four Mg\textsuperscript{2+} concentrations. At $[\mathrm{Mg}^{2+}]_o = 1\,\mathrm{mM}$
    (physiological), the channel is $>$90\% blocked at $V = -70\mV$ but $>$80\%
    unblocked at $V = 0\mV$. \textbf{(B)} I-V relationship: the resulting
    N-shaped I-V curve is diagnostic of NMDA receptor-mediated current and
    is qualitatively absent at zero Mg\textsuperscript{2+}.}
  \label{fig:nmda}
\end{figure}

The Mg\textsuperscript{2+} block (\Cref{fig:nmda}) is perhaps the most
biologically important non-linearity in the entire nervous system. It renders
NMDA a molecular ``AND gate'': simultaneous presynaptic glutamate release
\emph{and} postsynaptic depolarisation are required to produce significant
Ca\textsuperscript{2+} influx. This coincidence-detection property is
considered the synaptic implementation of Hebb's postulate \citep{Hebb1949}
and is mechanistically required for associative LTP \citep{Bliss1993}.
The N-shaped I-V curve in \Cref{fig:nmda}B carries intrinsic bistability:
at intermediate voltages, negative-slope conductance regions can sustain
two stable membrane potential states, hypothesised to underlie working
memory \emph{up-states} \citep{Wang1999}.

\subsection{Figure 7 --- STDP Weight Dynamics}

\begin{figure}[H]
  \centering
  \includegraphics[width=\columnwidth]{fig_07_stdp.png}
  \caption{\textbf{Spike-timing dependent plasticity (STDP).}
    \textbf{(A)} Theoretical STDP learning window: LTP (red) for
    $\Delta t > 0$ (post follows pre), LTD (blue) for $\Delta t < 0$.
    The slight asymmetry $A_-/A_+ = 1.05$ ensures a net depression at
    equal rates, critical for weight stabilisation.
    \textbf{(B)} Single-synapse weight evolution under 20\Hz Poisson
    pre- and postsynaptic activity. The weight converges to $w^* \approx 0.50$.
    \textbf{(C)} Steady-state weight distribution across 50 independent
    synapses (median $\approx 0.48$): broad unimodal, consistent with
    the multiplicative STDP prediction of \citet{Song2000}.}
  \label{fig:stdp}
\end{figure}

The STDP learning window (\Cref{fig:stdp}A) is reproduced from the
seminal experimental results of \citet{Bi1998} in hippocampal cultures.
Critically, the temporal asymmetry encodes causal relationships: only
when a presynaptic neuron fires \emph{before} a postsynaptic neuron --- consistent
with the former \emph{causing} the latter --- is the synapse potentiated.
This implements a temporal credit-assignment rule that can, in principle,
learn sequences and associations \citep{Abbott2000}.

The convergence of weights to an intermediate value (\Cref{fig:stdp}B--C)
under balanced random firing reflects the competition between potentiation
and depression. The slight depression bias ($A_- > A_+$) stabilises the
weight by preventing runaway potentiation under the Poisson regime. In a
structured input regime (e.g., correlated pre/post activity), the STDP rule
would drive weights to a bimodal ``winner-take-all'' distribution, selecting
the most informative synaptic inputs \citep{Song2000}.

\subsection{Figure 8 --- E/I Network Dynamics}

\begin{figure}[H]
  \centering
  \includegraphics[width=\columnwidth]{fig_08_network_dynamics.png}
  \caption{\textbf{Recurrent E/I network dynamics.}
    \textbf{(A)} Raster plot: each row is a neuron, each dot is a spike.
    Red: excitatory (100); blue: inhibitory (25). Dashed line separates
    populations. \textbf{(B)} Population firing rates (10-ms Gaussian smoothing):
    E $\approx 21\Hz$, I $\approx 26\Hz$. \textbf{(C)} LFP proxy (mean E-cell
    voltage), showing sub-threshold fluctuations. \textbf{(D)} E/I balance:
    red/blue fills indicate E- or I-dominant epochs, revealing dynamic competition
    between populations.}
  \label{fig:network}
\end{figure}

The E/I network simulation (\Cref{fig:network}) demonstrates that recurrent
excitation and inhibition together generate sustained asynchronous activity
without external drive beyond the tonic Poisson input. The inhibitory
population fires slightly faster (26\Hz vs 21\Hz) than excitatory neurons,
consistent with the inhibitory stabilisation mechanism \citep{Tsodyks1997}:
strong feedback inhibition prevents excitatory runaway and constrains the
network to a stable operating regime.

The population firing rates exhibit fluctuations on a timescale of $\sim$20\ms
that reflect the recruitment and suppression of activity cascades through the
recurrent network. The LFP proxy captures the aggregate subthreshold activity,
reflecting the dominant timescale of synaptic currents.

\subsection{Figure 9 --- ISI Statistics and Firing Irregularity}

\begin{figure}[H]
  \centering
  \includegraphics[width=\columnwidth]{fig_09_isi_analysis.png}
  \caption{\textbf{ISI statistics.}
    \textbf{(A)} Log-scale ISI histogram for all E-cells pooled. The
    right-skewed distribution with a mode at $\sim$25\ms (40\Hz) is consistent
    with a mostly irregular but rate-dominated spiking regime.
    \textbf{(B)} Poincaré return map (ISI$_{n+1}$ vs ISI$_n$) for the
    most active E-cell. Diffuse scatter indicates moderate irregularity.
    \textbf{(C)} CV-ISI distribution across 100 E-neurons. Median CV $= 0.18$
    indicates a \emph{synchronous-regular} (SR) state rather than true AI
    ($\mathrm{CV} \approx 1$), reflecting the predominantly feed-forward
    Poisson drive.}
  \label{fig:isi}
\end{figure}

The ISI analysis (\Cref{fig:isi}) reveals the network is operating in a
synchronous-regular (SR) state (CV-ISI $\approx 0.18$), rather than the
biologically desirable AI state (CV-ISI $\approx 1$). This is expected: our
simplified network of 125 neurons with relatively strong Poisson drive produces
predominantly driven firing rather than self-sustained irregular activity. The
Poincaré return map is tightly clustered near the diagonal, confirming low
trial-to-trial variability.

The critical parameter for transitioning to the AI state is the ratio of
inhibitory to excitatory strength $g = g_{IE}/g_{EE}$ and the external firing
rate \citep{Brunel2000}. For the AI regime, the effective synaptic noise from
thousands of recurrent inputs must dominate the mean input — a condition
satisfied only in much larger ($N \sim 10^4$) networks or with sufficiently
decorrelated inputs.

\subsection{Figure 10 --- LFP Power Spectrum}

\begin{figure}[H]
  \centering
  \includegraphics[width=\columnwidth]{fig_10_lfp_psd.png}
  \caption{\textbf{LFP power spectral density.}
    \textbf{(A)} Welch PSD of the LFP proxy. Frequency bands are
    colour-coded: $\delta$ (0.5--4\Hz, green), $\theta$ (4--8\Hz, purple),
    $\alpha$ (8--13\Hz, red), $\beta$ (13--30\Hz, blue), $\gamma$ (30--100\Hz,
    orange). The power is concentrated in the $\gamma$ band, consistent with the
    $\sim$25\Hz dominant network oscillation emerging from the E/I feedback loop.
    \textbf{(B)} Integrated band powers: $\gamma$ dominates by several orders
    of magnitude over lower frequency bands.}
  \label{fig:psd}
\end{figure}

The LFP power spectrum (\Cref{fig:psd}) is dominated by $\gamma$-band power
($>$30\Hz), consistent with the $\sim$20\Hz mean E-cell firing rate and the
E/I feedback oscillation frequency. In real cortex, gamma oscillations arise
through the ING (Interneuron Network Gamma) and PING (Pyramidal Interneuron
Network Gamma) mechanisms \citep{Whittington2000}, both of which involve the
interplay of fast excitation (AMPA) and delayed inhibition (GABA\textsubscript{A})
at the network level. The dominant frequency of PING oscillations is approximately
$f \approx 1 / (t_{\mathrm{delay}} + \tau_{\mathrm{GABA}})$, consistent with
our result.

%% ──────────────────────────────────────────────────────────
\section{Discussion and Scientific Significance}
\label{sec:discussion}
%% ──────────────────────────────────────────────────────────

\subsection{Model Hierarchy and Computational Trade-offs}

The three neuron models implemented span a fundamental trade-off between
biophysical detail and computational efficiency (\Cref{tab:model_comparison}).
The HH model, with its four-dimensional ODE system and voltage-dependent
conductances, is the gold standard for membrane dynamics, but requires RK4
integration at $\Delta t \leq 0.025\ms$ (40 kHz) to resolve the fast
Na\textsuperscript{+} activation. The LIF reduces the membrane to a
one-dimensional linear ODE with a reset, enabling large-scale network
simulation ($N \sim 10^5$ on a single CPU) while retaining the essential
integrate-and-fire computation. The AdEx bridges these extremes: its
exponential term captures realistic spike initiation, and the adaptation
variable $w$ recovers the full diversity of cortical firing patterns at
only $2\times$ the computational cost of LIF.

\begin{table}[H]
\centering
\small
\caption{Neuron model comparison.}
\label{tab:model_comparison}
\begin{tabularx}{\columnwidth}{lccXc}
\toprule
\textbf{Model} & \textbf{Dim.} & \boldmath$\Delta t$ & \textbf{Properties captured} & \textbf{Biophys.} \\
\midrule
HH  & 4 & 0.025\ms & AP shape, channel kinetics, refractory & **** \\
LIF & 1 & 0.1\ms & Subthreshold integration, rate coding & ** \\
AdEx & 2 & 0.1\ms & Adaptation, firing patterns, EIF & *** \\
\bottomrule
\end{tabularx}
\end{table}

\subsection{The Role of Inhibition in Network Stability}

The E/I network results highlight a fundamental principle: \emph{inhibition
is not merely the absence of excitation, but an active computational mechanism}.
Without recurrent inhibition, the network would undergo synchronised epileptic
bursting (as readily demonstrated by setting $g_{IE} = 0$). With appropriate
inhibition, three outcomes are possible \citep{Brunel2000}:

\begin{enumerate}[leftmargin=1.2em, itemsep=2pt]
  \item \textbf{AI state} ($g \geq 5$, low external rate): individual
    neurons fire irregularly due to the near-cancellation of excitation and
    inhibition, producing high CV-ISI $\approx 1$.
  \item \textbf{SR state} ($g < 3$, high external rate): network is
    driven primarily by external input and fires regularly.
  \item \textbf{Oscillatory state} (intermediate $g$): E/I loop generates
    gamma oscillations (PING mechanism).
\end{enumerate}

Our simulation operates in the SR state at the parameters chosen, reflecting
the relatively strong Poisson drive. The transition to AI requires the effective
synaptic noise --- the standard deviation of recurrent input currents --- to
exceed the distance between resting potential and threshold.

\subsection{STDP and the Emergence of Selectivity}

The STDP simulations demonstrate that competitive Hebbian learning under
random input drives weights to an intermediate stable value. This changes
dramatically under structured input: if a subset of presynaptic neurons
fire in a correlated temporal pattern (e.g., a sequence), STDP selectively
potentiates those synapses while depressing others \citep{Song2000}. Over
time, this produces neurons that respond preferentially to specific temporal
patterns --- a plausible mechanism for \emph{feature selectivity} in sensory
cortex.

The key insight is that the asymmetric temporal window encodes
\emph{causality}: synapses from neurons that reliably predict the postsynaptic
neuron are strengthened, while non-predictive synapses are weakened. This is
a form of \emph{temporal contrastive Hebbian learning} \citep{Abbott2000}.

\subsection{NMDA as the Molecular Substrate of Learning}

The voltage-dependent Mg\textsuperscript{2+} block (\Cref{fig:nmda})
is mechanistically indispensable for associative plasticity: it ensures that
only synapses active during postsynaptic firing (i.e., those that contributed
to the AP) are eligible for LTP. Blocking NMDA receptors with APV abolishes
LTP induction in hippocampal slices \citep{Bliss1993}, and knockin mice
with inactivated NMDA receptors in CA1 pyramidal neurons are severely impaired
on spatial learning tasks \citep{Tsien1996}. Our implementation faithfully
reproduces the N-shaped I-V curve (\Cref{fig:nmda}B) and can therefore
serve as the synaptic plasticity module in more complex learning simulations.

\subsection{Limitations and Extensions}

Several biophysically important features are not included in the current
implementation:

\begin{itemize}[leftmargin=1.2em, itemsep=2pt]
  \item \textbf{Dendritic computation}: real neurons perform nonlinear
    dendritic integration. Multi-compartment models \citep{Mainen1996}
    would capture synaptic location effects.
  \item \textbf{Short-term plasticity}: synaptic facilitation and depression
    on timescales of 100--1000\ms \citep{Tsodyks1997} are absent.
  \item \textbf{Calcium dynamics}: a full Ca\textsuperscript{2+} model
    \citep{Bhatt2009} would tie NMDA to LTP magnitude.
  \item \textbf{Neuromodulation}: dopamine, acetylcholine, and norepinephrine
    modulate both intrinsic and synaptic properties. Their inclusion is
    necessary for modelling arousal, attention, and reward.
\end{itemize}

%% ──────────────────────────────────────────────────────────
\section{Software Architecture}
\label{sec:software}
%% ──────────────────────────────────────────────────────────

\subsection{Project Structure}

\begin{lstlisting}[caption={Project directory structure.},
                   language=bash, basicstyle=\ttfamily\scriptsize]
compneuro/
 neurons/
   base_neuron.py    # Abstract base class
   hodgkin_huxley.py # HH model + RK4 + f-I
   integrate_fire.py # LIF, EIF, AdEx presets
 synapses/
   synapse.py        # AMPA, NMDA, GABA-A/B kinetics
   plasticity.py     # STDP, BCM, Oja, Triplet rules
 networks/
   network.py        # Sparse random E/I network
 analysis/
   spike_analysis.py # ISI, PSTH, autocorr, PSD
   visualisation.py  # 11 publication figures
 database/
   models.py         # SQLAlchemy ORM schema
   db.py             # Engine, sessions, repository
 generate_figures.py # Master simulation script
 requirements.txt    # Python dependencies
\end{lstlisting}

\subsection{Database Schema}

Results are persisted to a SQLite database via SQLAlchemy ORM, enabling
reproducible query of simulation metadata:

\begin{lstlisting}[caption={Querying simulation results via SQLAlchemy.}]
from database.db import build_engine, SimulationRepository

engine = build_engine("compneuro.db")
repo   = SimulationRepository(engine)

# Retrieve all HH simulation runs
hh_runs = repo.list_runs(model_type="HodgkinHuxley")

# Get spike train statistics for run 1
summary = repo.run_summary(run_id=1)
print(f"Mean rate: {summary['mean_rate_hz']:.1f} Hz")
print(f"CV-ISI   : {summary['mean_cv_isi']:.3f}")
\end{lstlisting}

\subsection{Reproducing the Results}

\begin{lstlisting}[caption={Running the simulation suite.}, language=bash]
# Install dependencies
pip install numpy scipy matplotlib sqlalchemy

# Clone repository
git clone https://github.com/[user]/compneuro
cd compneuro

# Run all simulations and generate figures
python generate_figures.py
# Output: figures/fig_01_*.png ... fig_11_*.png
#         compneuro.db
\end{lstlisting}

\subsection{Key Implementation Decisions}

\begin{notebox}[title=Engineering Notes]
\textbf{RK4 vs Euler integration.} The Euler method introduces $\mathcal{O}(\Delta t)$
local truncation error, which accumulates dangerously near the spike threshold.
RK4 reduces this to $\mathcal{O}(\Delta t^4)$, allowing stable integration
at $\Delta t = 0.025\ms$ for the HH model. For LIF and AdEx, Euler at
$\Delta t = 0.1\ms$ is adequate.

\textbf{AdEx overflow prevention.} The exponential term
$\gL\dT e^{(V-V_T)/\dT}$ can overflow double precision if $V > +30\mV$.
The fix is to immediately reset $V \leftarrow \Vreset$ upon threshold crossing
within the same timestep, preventing the voltage from persisting at $V_{\mathrm{peak}}$.

\textbf{SQLAlchemy WAL mode.} SQLite's write-ahead logging (WAL) is
enabled via a connection pragma, enabling concurrent read/write during
long simulation runs.
\end{notebox}

%% ──────────────────────────────────────────────────────────
\section{Conclusions}
\label{sec:conclusions}
%% ──────────────────────────────────────────────────────────

We have implemented and validated a complete biophysical simulation
hierarchy spanning four levels of neural organisation. The key scientific
conclusions are:

\begin{enumerate}[leftmargin=1.4em, itemsep=3pt]
  \item The Hodgkin--Huxley model faithfully reproduces AP shape,
    ion-current anatomy, and f-I curve rheobase from the original
    1952 parameterisation, confirming correctness of the RK4
    implementation.

  \item The AdEx model reproduces five distinct cortical firing-pattern
    archetypes (RS, IB, CH, FS, LTS) by varying only three parameters
    $(a, b, \tauw)$, demonstrating the model's utility as a flexible
    yet tractable cell-class descriptor.

  \item The kinetic synapse models reproduce the published time constants
    and I-V characteristics of AMPA, NMDA, GABA\textsubscript{A}, and
    GABA\textsubscript{B} receptors. The NMDA Mg\textsuperscript{2+}
    block is an essential non-linearity for both coincidence detection
    and associative learning.

  \item STDP with multiplicative weight bounds and a slight depression
    bias ($A_-/A_+ = 1.05$) produces a unimodal steady-state weight
    distribution, confirming stability under unstructured Poisson input.

  \item The 125-neuron sparse random E/I LIF network exhibits
    self-sustained activity at biologically plausible rates ($\sim$20\Hz)
    with gamma-band dominant LFP, demonstrating the emergence of network
    oscillations from purely local E/I interactions.
\end{enumerate}

The codebase is designed as a modular, extensible research platform.
Immediate extensions include multi-compartment dendritic integration,
short-term synaptic plasticity, and calcium-dependent LTP rules,
all of which can be integrated without modification of the existing
model hierarchy.

%% ──────────────────────────────────────────────────────────
\section*{Acknowledgements}
%% ──────────────────────────────────────────────────────────

The authors acknowledge the foundational contributions of
\citeauthor{HodgkinHuxley1952}, \citeauthor{Brette2005},
\citeauthor{Destexhe1994}, \citeauthor{Bi1998}, and \citeauthor{Brunel2000},
whose theoretical frameworks are implemented in this work.
Figures were generated using \textsc{Matplotlib} \citep{Hunter2007}.
Database management uses SQLAlchemy \citep{Bayer2012}.
Numerical computations used NumPy \citep{Harris2020} and SciPy \citep{Virtanen2020}.

%% ──────────────────────────────────────────────────────────
\bibliography{references}
%% ──────────────────────────────────────────────────────────

\end{document}
